\documentclass[11pt]{article}
\usepackage[margin=1in]{geometry}
\usepackage{amsmath,amssymb,amsthm,mathtools,mathrsfs}
\usepackage{graphicx}
\usepackage{microtype}
\usepackage[colorlinks=true,linkcolor=blue,citecolor=blue,urlcolor=blue]{hyperref}

\newtheorem{theorem}{Theorem}
\newtheorem{lemma}{Lemma}
\newcommand{\D}{\mathbb D}
\newcommand{\Berg}{\mathscr B}
\newcommand{\HH}{\mathscr H}
\newcommand{\T}{\mathsf T}
\newcommand{\e}{\mathrm e}
\newcommand{\ii}{\mathrm i}
\newcommand{\Span}{\operatorname{span}}
\newcommand{\tr}{\operatorname{tr}}

\title{A Strictly Positive-Field Counterexample to the\\
Boundary--Bergman Equality Question}
\author{Automated mathematics run \texttt{fa\_banach\_001}, lane 10}
\date{13 August 2026}

\begin{document}
\maketitle

\begin{center}
\fbox{\parbox{0.91\textwidth}{\textbf{Status.}
Candidate full counterexample, likely valid, pending human review.  A smooth
polynomial magnetic field satisfying $B>0$ makes the first inequality of
Fournais--Frank--Goffeng--Kachmar--Sundqvist Proposition~5.2 strict.}}
\end{center}

\section{Source question}

Fournais, Frank, Goffeng, Kachmar, and Sundqvist study magnetic Pauli
operators on planar domains in \cite{FFGKS}.  Their Proposition~5.2 gives, in
particular,
\begin{equation}\label{source-ineq}
 N(\HH^{\mathcal O}_{\mathbf A,g},0)
 \geq
 \dim\!\left(E_V\cap\gamma_0(\Berg(\Omega))\right),
 \qquad V=g-A_\tau .
\end{equation}
Remark~5.3 on PDF page 18 asks whether the inequalities in that proposition
are equalities when $B\geq0$.  The latest source version renumbers the same
passage as Proposition~6.3 and Remark~6.4.

\begin{figure}[ht]
  \centering
  \includegraphics[width=\textwidth]{figures/open_problem_crop.png}
  \caption{Full-width source crop from arXiv:2307.16079v1, PDF page 18,
  containing Remark~5.3 and the exact equality question.}
\end{figure}

We disprove the proposed first equality with a field that is strictly
positive on the closed disc.  This suffices to answer the general question
negatively.

\section{Definitions and proof intuition}

Let $\D=\{z\in\mathbb C:|z|<1\}$ and impose Neumann data $g=0$.  Following
the source, choose a real function $\phi$ vanishing on $\partial\D$ and put
\[
 \mathbf A=(-\partial_y\phi,\partial_x\phi),
 \qquad B=\operatorname{curl}\mathbf A=\Delta\phi.
\]
The source orients the boundary tangent counterclockwise (and the normal
inward), so
\[
 A_\tau=\mathbf A\cdot\tau=\partial_r\phi\big|_{r=1}.
\]
On holomorphic functions the compressed quadratic form is
\begin{equation}\label{qform}
 q^{\mathcal O}_{\mathbf A,0}(v)
 =\int_0^{2\pi}
 \left((-\ii\partial_\theta-A_\tau(\theta))v(\e^{\ii\theta})\right)
 \overline{v(\e^{\ii\theta})}\,d\theta.
\end{equation}

\paragraph{Proof intuition.}
We prescribe $A_\tau$ to have only Fourier frequencies $0$ and $\pm3$ and
realize it by a polynomial $\phi$ whose Laplacian remains positive.  Formula
\eqref{qform} then splits into three blocks according to Fourier index modulo
three.  Two blocks contain the plainly negative modes $1,z$, while a
$2\times2$ determinant forces a negative direction in the third block.
The comparison space in \eqref{source-ineq} is nevertheless exactly
two-dimensional: the explicit boundary-Dirac eigenfunctions admit a
Wiener--Hopf factorization, reducing holomorphicity to the preimage of
$\Span\{1,z\}$ under an invertible coanalytic Toeplitz operator.

\section{Counterexample theorem}

\begin{theorem}\label{main-thm}
Let
\[
 a=\frac{399}{200},\qquad b=\frac6{25},
\]
and define on the unit disc
\begin{equation}\label{phi-def}
 \phi(r,\theta)
 =\frac a2(r^2-1)-b(r^3-r^5)\cos(3\theta),
 \qquad \mathbf A=(-\phi_y,\phi_x).
\end{equation}
Then $\phi$ and $\mathbf A$ are smooth polynomials, $\phi|_{\partial\D}=0$,
and
\[
 B=\operatorname{curl}\mathbf A\geq\frac3{20}>0
 \quad\hbox{on }\overline\D.
\]
For Neumann data $g=0$, the first inequality in Proposition~5.2 of
\cite{FFGKS} is strict:
\begin{equation}\label{strict}
 N(\HH^{\mathcal O}_{\mathbf A,0},0)
 \geq3>2
 =\dim\!\left(E_{-A_\tau}\cap\gamma_0(\Berg(\D))\right).
\end{equation}
In particular, nonnegativity of the magnetic field does not imply equality.
\end{theorem}

\begin{proof}
We divide the proof into the field construction, the negative-subspace
calculation, and the exact comparison-space calculation.

\paragraph{1. The field is polynomial and strictly positive.}
Writing
\[
 P_3(x,y)=x^3-3xy^2=r^3\cos(3\theta),
\]
formula \eqref{phi-def} becomes
\[
 \phi(x,y)=\frac a2(x^2+y^2-1)
 -b\bigl(1-x^2-y^2\bigr)P_3(x,y),
\]
so it is a polynomial and vanishes on the circle.  Direct differentiation
gives
\begin{equation}\label{boundary-and-field}
 A_\tau(\theta)=\partial_r\phi(1,\theta)
 =a+2b\cos(3\theta),
 \qquad
 B(r,\theta)=\Delta\phi=2a+16br^3\cos(3\theta).
\end{equation}
Since $0\leq r\leq1$,
\[
 B(r,\theta)\geq2a-16b
 =\frac{399}{100}-\frac{96}{25}=\frac3{20}>0.
\]

\paragraph{2. The compressed form has at least three negative eigenvalues.}
For a holomorphic polynomial $v(z)=\sum_{n\geq0}c_nz^n$, Fourier
orthogonality and \eqref{boundary-and-field} give
\begin{equation}\label{fourier-form}
 \frac1{2\pi}q^{\mathcal O}_{\mathbf A,0}(v)
 =\sum_{n\geq0}(n-a)|c_n|^2
 -b\sum_{n\geq0}
 \left(c_{n+3}\overline{c_n}+c_n\overline{c_{n+3}}\right),
\end{equation}
where coefficients outside the polynomial's support are zero.  Thus the form
is block diagonal according to $n$ modulo three.

The directions $1$ and $z$ belong to two distinct residue blocks and have
form values $-2\pi a$ and $2\pi(1-a)$, both negative.  In the third residue
block, the matrix of $q^{\mathcal O}_{\mathbf A,0}/(2\pi)$ on
$\Span\{z^2,z^5\}$ is
\[
 M=\begin{pmatrix}2-a&-b\\-b&5-a\end{pmatrix}.
\]
Its trace is positive, whereas
\begin{align*}
 \det M
 &=(2-a)(5-a)-b^2\\
 &=\frac1{200}\frac{601}{200}-\left(\frac6{25}\right)^2
 =-\frac{1703}{40000}<0.
\end{align*}
Hence $M$ has one negative eigenvalue.  Choose a corresponding negative
vector $w\in\Span\{z^2,z^5\}$.  The three vectors $1,z,w$ lie in distinct
residue blocks, so they are mutually orthogonal for the quadratic form and
their span is negative definite.  The min--max principle (equivalently,
Glazman's lemma) proves
\[
 N(\HH^{\mathcal O}_{\mathbf A,0},0)\geq3.
\]

\paragraph{3. The boundary comparison space has dimension two.}
Set $V=-A_\tau=-a-2b\cos(3\theta)$.  The explicit spectral formula for
$\mathcal D_V=-\ii\partial_\theta+V$ in \cite[Proposition~3.1]{FFGKS}
gives eigenvalues $m-a$, $m\in\mathbb Z$, and corresponding eigenfunctions,
up to normalization,
\begin{equation}\label{eigenfunctions}
 f_m(z)=z^mE(z),\qquad
 E(z)=\exp\!\left(\delta(z^3-z^{-3})\right),
 \qquad \delta=\frac b3,quad |z|=1.
\end{equation}
Indeed, the source formula uses
$\int_0^\theta V=-a\theta-(2b/3)\sin(3\theta)$.
Because $1<a<2$, the negative eigenvalues are precisely those with $m\leq1$.
Thus
\begin{equation}\label{EV-form}
 E_V=E\,\overline{\Span\{z^m:m\leq1\}}^{\,L^2(\partial\D)}.
\end{equation}

Let $H^2$ denote the usual Hardy space of nonnegative boundary Fourier modes.
The source proves that
$E_V\cap\gamma_0(\Berg(\D))\subset H^{1/2}(\partial\D)$; hence every vector
in the intersection is in $H^2$.  Conversely an $H^2$ boundary function is
the trace of a holomorphic Hardy function and therefore of a Bergman
function.  It is consequently enough to calculate $E_V\cap H^2$.

Factor
\[
 E=E_+E_-,\qquad
 E_+(z)=\exp(\delta z^3),\qquad
 E_-(z)=\exp(-\delta z^{-3}).
\]
Suppose $f=Eh\in H^2$, where $h$ has Fourier support in $\{m:m\leq1\}$ as
in \eqref{EV-form}, and put $G=E_+^{-1}f\in H^2$.  Then
\[
 h=E_-^{-1}G=\overline{\alpha}\,G,
 \qquad \alpha(z)=\exp(\delta z^3)=E_+(z)quad(|z|=1).
\]
Writing $P_+$ for the Hardy projection, the support condition on $h$ is
equivalent to
\begin{equation}\label{toeplitz-condition}
 \T_{\overline\alpha}G
 :=P_+(\overline\alpha G)\in\Span\{1,z\}.
\end{equation}
The functions $\alpha$ and $\alpha^{-1}$ are bounded analytic functions on
$\D$.  Coanalytic Toeplitz operators multiply in the same order, and hence
\[
 \T_{\overline\alpha}\T_{\overline{\alpha^{-1}}}
 =\T_{\overline{\alpha\alpha^{-1}}}=I,
 \qquad
 \T_{\overline{\alpha^{-1}}}\T_{\overline\alpha}=I.
\]
Therefore $\T_{\overline\alpha}$ is invertible.  The solution space of
\eqref{toeplitz-condition} is the inverse image of the two-dimensional space
$\Span\{1,z\}$ and consequently has dimension exactly two.  There is no
regularity deficit: for $G=1,z$, the vectors
\[
 f=E_+, zE_+
\]
are smooth holomorphic traces, while
$h=E_-^{-1},zE_-^{-1}$ has Fourier support at most one.  We conclude that
\[
 \dim\!\left(E_V\cap\gamma_0(\Berg(\D))\right)=2.
\]
Together with the preceding lower bound this proves \eqref{strict}.
\end{proof}

\section{A parameter family and robustness}

The rational constants in Theorem~\ref{main-thm} are only a convenient exact
choice.  The proof gives the following family.

\begin{lemma}\label{family}
Use \eqref{phi-def} with real parameters satisfying
\[
 1<a<2,\qquad b\neq0,\qquad a\geq8|b|,
 \qquad (2-a)(5-a)<b^2.
\]
Then $B\geq0$, the comparison dimension is two, and the compressed form has
at least three negative eigenvalues.  If $a>8|b|$, then $B>0$ on
$\overline\D$.
\end{lemma}

\begin{proof}
Equations \eqref{boundary-and-field}--\eqref{fourier-form} remain valid.  The
condition $a\geq8|b|$ is exactly $2a-16|b|\geq0$.  The determinant condition
forces a negative direction in the $z^2,z^5$ block, while $1<a<2$ gives the
same two negative low modes and the same cutoff $m\leq1$ in the Dirac
spectrum.  The Toeplitz argument is unchanged.
\end{proof}

For each sufficiently small fixed $b\neq0$, taking $a<2$ sufficiently close
to two satisfies the determinant condition; the positivity condition is then
also satisfied.  The strict inequalities in Theorem~\ref{main-thm} show that
its example lies inside an open set of positive-field counterexamples.

\section{Verification, scope, and novelty}

The bundled script \texttt{code/exact\_symbolic\_check.py} verifies with exact
symbolic and rational arithmetic the polynomial Laplacian, the boundary
tangential potential, the lower bound $B\geq3/20$, and the determinant
$-1703/40000$.  The Hardy/Toeplitz calculation is an analytic proof and is
not inferred from a truncation.

Only the first equality proposed in Remark~5.3 is refuted.  This is enough to
answer the general question negatively.  We do not claim a characterization
for equality in the second, antiholomorphic inequality.

A bounded novelty search on 13 August 2026 covered the local run indexes and
attempts, exact-title and exact-phrase web/arXiv queries, the arXiv id, the
remark numbers, all five authors, and combinations of the core terms
``nonnegative magnetic field'', ``Bergman compression'', and ``equality''.
It found the source preprint and its 2025 journal publication, but no later
claimed resolution or matching counterexample.  The latest source TeX still
contains the question.  This supports moderate, not definitive, novelty
confidence.

Human review should focus on the boundary orientation sign, the cutoff
$m\leq1$ in \eqref{EV-form}, and the equivalence
\eqref{toeplitz-condition}.  No unproved dependency is currently known.

\begin{thebibliography}{9}
\bibitem{FFGKS}
S.~Fournais, R.~L. Frank, M.~Goffeng, A.~Kachmar, and M.~Sundqvist,
\emph{Counting Negative Eigenvalues for the Magnetic Pauli Operator},
arXiv:2307.16079v1 (2023); Duke Math. J. \textbf{174} (2025), 313--353.
\end{thebibliography}

\end{document}
